\chapter*{ПРИЛОЖЕНИЕ А}
\addcontentsline{toc}{chapter}{ПРИЛОЖЕНИЕ А}

\noindent\textbf{Структура репозитория и запуск}

{\footnotesize\begin{verbatim}
1-dataset/          датасет (ЛР1): dataset.json, xlsx-копия
2-metrics/          меры близости (ЛР2): metrics.py, out/ (матрицы,
                    тепловые карты, matrices.html), docs/
3-likesdislikes/    рекомендательная система (ЛР3)
  run.py            запуск интерфейса, http://127.0.0.1:8003/
  recsys/           domain, similarity, recommend, service, web
  experiments/      evaluate.py, plots.py -> out/
  tests/            unittest (25 тестов)
  tools/            method_length.py - правило "<= 5 строк"
4-search/           поиск и объединённая система (ЛР4, РК1)
  run.py            запуск интерфейса, http://127.0.0.1:8004/
  bootstrap.py      подключение пакета recsys из ЛР3
  search/           фильтры, запрос, полезность, ослабление
  unified/          сессии, история, команды, API, интерфейс
  experiments/      compare_fallback.py -> out/
  tests/            unittest (13 тестов)
report/             этот отчёт (LaTeX)
\end{verbatim}}

\noindent Требования: Python 3.10 и выше; сторонние пакеты для работы систем не нужны. Для проверочных расчётов и ЛР2 нужны \code{numpy}, \code{matplotlib}, \code{openpyxl}.

{\footnotesize\begin{verbatim}
cd 3-likesdislikes && python run.py
cd 4-search && python run.py
python -m unittest discover -s tests -t .
python experiments/evaluate.py && python experiments/plots.py
\end{verbatim}}

\noindent Рисунки отчёта собираются скриптом \code{report/tools/build\_figures.py}: диаграммы из \code{report/diags/*.dot} (Graphviz), графики сравнений подходов и тепловые карты ЛР2 копируются в \code{report/img/}.
